% Setting up the document class and essential packages
\documentclass[a4paper,12pt]{article}
\usepackage{amsmath, amsthm, amssymb,color,braket}
\usepackage{geometry}
\geometry{margin=0.5in} % Set 1.5-inch margins on all sides

\usepackage{parskip}
\renewcommand{\familydefault}{\sfdefault}


\newtheoremstyle{mystyle}%                % Name
{}%                                     % Space above
{}%                                     % Space below
{\itshape}%                             % Body font
{}%                                     % Indent amount
{\bfseries}%                            % Theorem head font
{.}%                                    % Punctuation after theorem head
{ }%                                    % Space after theorem head, ' ', or \newline
{}%     


\theoremstyle{mystyle}
% Configuring theorem-like environments
\newtheorem{theorem}{Theorem}[section]
\newtheorem{lemma}[theorem]{Lemma}
\newtheorem{corollary}[theorem]{Corollary}

\theoremstyle{mystyle}
\newtheorem{proposition}{Proposition}[section]
\newtheorem{definition}{\color{blue} Definition}[section]
\newtheorem{remark}{\color{red} Remark}[section]
\newtheorem{example}{Example}[section]
\title{Slater Neural Quantum States}
% Setting up section numbering
\numberwithin{equation}{section}
\begin{document}
	\maketitle
	We give a brief introduction on how statistic reconfiguration (SR) is implemented in optimization of NQS.Take any random initial state $\ket{\psi_{\tau=0}}$,if we keep projecting the state as follows
	$$\frac{d}{d\tau}\ket{\psi}=(\Lambda-H)\ket{\psi}$$
	Then the solution is $\ket{\psi_\tau}=e^{\Lambda-H}\ket{\psi}$.From the solution we see that the state will eventually converges to ground state.When updating the NQS we want the learning step to be "similar" to imaginary time evolution,that is
	$$\delta_\theta \partial_\theta \ket{\psi}\sim \eta(\Lambda-H)\ket{\psi}$$
\end{document}